% =====================================================================
%  Chương 8 — Tác tử học tăng cường và tác tử dựa trên luật (Agent 02)
% =====================================================================
\chapter{Tác tử học tăng cường và tác tử dựa trên luật}\label{chap:rl}

Chương này cô đọng bài tập Agent 02: hai cửa hàng cùng quyết định giữ giá hay giảm giá, một bên dùng Q-learning, một bên dùng
luật viết tay. Môi trường là bộ mô phỏng theo đề bài~\cite{ch8:lab02}, mọi quan sát đều là \textbf{dữ liệu tổng hợp}. Số liệu dưới đây được chép
nguyên từ lượt chạy của bài tập đó, không có thí nghiệm mới.

\section{Mô hình hóa}

Ở mỗi bước $t$, cửa hàng $i$ chọn đồng thời $a_i\in\{0,1\}$ (0: giữ giá, 1: giảm giá) sau khi quan sát $s_i=(d_t,c_i)$, với
$d_t\in\{0,1\}$ là nhu cầu và $c_i=a_{-i}(t-1)$ là hành động trước của đối thủ. Số khách $N_t$ là 10 hoặc 20; số đơn bán được
$S_i\sim\mathrm{Binomial}(N_t,p_i)$ với $p_i$ ở Hình~\ref{fig:ch8-loop}. Phần thưởng là lợi nhuận với $P=10$, $C=5$, $\delta=0{,}20$; mỗi episode có 20 bước.
Hai tác tử không trao đổi thông điệp mà ảnh hưởng nhau qua thị trường chung.

\hinhtikz{ch8_loop}{Một bước tương tác: hai tác tử chọn hành động đồng thời, thị trường tính xác suất chuyển đổi, doanh số và phần thưởng.}{ch8-loop}

Tác tử luật giữ giá khi nhu cầu thấp; khi nhu cầu cao, nó chỉ giảm giá nếu đối thủ vừa giảm. Tác tử RL giữ bảng $Q$ tám ô và cập nhật
theo Q-learning~\cite{ch8:watkins1992} với $\alpha=0{,}10$, $\gamma=0{,}90$, $\varepsilon_e=\max(0{,}02;\,0{,}20\cdot0{,}995^e)$:
\begin{equation}\label{eq:ch8-q}
  Q(s_t,a_t)\leftarrow Q(s_t,a_t)+\alpha\bigl[r_t+\gamma\max_{a'}Q(s_{t+1},a')-Q(s_t,a_t)\bigr].
\end{equation}

\maNguon{ch8_agents__QLearningAgent.learn}{Cập nhật Q-learning theo \eqref{eq:ch8-q}.}

Vì mô hình nhỏ, lợi nhuận kỳ vọng tính được chính xác (Bảng~\ref{tab:ch8-reward}). Với $\delta=20\%$, giữ giá tốt hơn ở ba trên bốn ô, nên không bao giờ giảm giá là
chính sách tối ưu với giá trị \SL{ch8_model/d20/optimal} mỗi bước, và tác tử luật vốn đã làm đúng như vậy. Với $\delta=5\%$, tối ưu chuyển thành luôn giảm giá
(\SL{ch8_model/d05/optimal}).

\bang{ch8_reward}{Lợi nhuận kỳ vọng một bước theo $(d,b)$; $b$ là hành động hiện tại của đối thủ.}{ch8-reward}

\section{Thí nghiệm A--E}

Sau huấn luyện 5.000 episode, bảng $Q$ đóng băng, $\varepsilon=0$, đánh giá \SL{ch8_a/steps} bước với seed riêng. Tác tử luật được đặt vào đúng ghế của RL trên cùng chuỗi ngẫu
nhiên để so cặp. Ngoài thưởng lấy mẫu, mỗi chính sách còn có giá trị chính xác, tính bằng lan truyền phân phối trạng thái.

\textbf{A, B.} Lượt tham chiếu đạt \SL{ch8_a/eval_mean} mỗi bước, tỉ lệ giảm giá \SL{ch8_a/eval_discPct}\%, trùng khít tác tử luật trên cùng chuỗi ngẫu nhiên; giá trị chính xác là \SL{ch8_a/exact}
(Hình~\ref{fig:ch8-q}). Thưởng huấn luyện tăng từ \SL{ch8_b/first500} (500 episode đầu) lên \SL{ch8_b/last500} (500 episode cuối). Đường học rất nhiễu vì nhiễu Binomial lớn hơn mức cải thiện do học
(Hình~\ref{fig:ch8-learning}, trái).

\hinhtikz{ch8_qbar3d}{Bảng $Q$ của lượt tham chiếu: ở cả bốn trạng thái, cột giữ giá cao hơn cột giảm giá.}{ch8-q}

\textbf{C, D.} Bảng~\ref{tab:ch8-sweeps} và Hình~\ref{fig:ch8-sweeps}: $\alpha=0{,}50$ cho thưởng thấp nhất và chính sách đổi \SL{ch8_c/0.50/switches} lần, vì mỗi mẫu nhiễu chiếm một nửa
giá trị cũ; $\alpha=0{,}01$ ổn định nhất nhưng thưởng huấn luyện chỉ \SL{ch8_c/0.01/train_mean}. Không khám phá là cấu hình tệ nhất (\SL{ch8_d/khong/n_opt}/10 seed tối ưu), vì hành động đầu tiên có thưởng dương
được giữ mãi; lịch thấp tốt nhất vì bảng chỉ có tám ô.

\bang{ch8_sweeps}{Thí nghiệm C, D, E: 10 seed, đánh giá tham lam; ``Đổi CS'' là số lần chính sách đổi khi huấn luyện.}{ch8-sweeps}
\hinhpdf{ch8_sweeps}{1}{Thưởng đánh giá theo $\alpha$ và theo lịch khám phá; chấm đen là từng seed, nét đỏ là mức tối ưu 35.}{ch8-sweeps}

\textbf{E.} Trên 10 seed, RL đạt \SL{ch8_e/agg/eval_mean} $\pm$ \SL{ch8_e/agg/eval_sd}, luật đạt \SL{ch8_e/agg/rule_mean}; chênh lệch cặp \SL{ch8_e/agg/diff_mean}, khoảng tin cậy 95\%
$[\SL{ch8_e/agg/ci_lo};\ \SL{ch8_e/agg/ci_hi}]$ chứa 0, nên chưa đủ bằng chứng nói bên nào hơn. \SL{ch8_e/agg/n_opt}/\SL{ch8_e/agg/n} seed đạt tối ưu; một seed học hỏng kéo trung bình xuống
(Hình~\ref{fig:ch8-learning}, phải).

\hinhpdf{ch8_learning}{1}{Trái: đường học lượt tham chiếu, trung bình trượt 100 episode. Phải: thưởng đánh giá trên mười seed.}{ch8-learning}

\section{Trả lời mười câu hỏi}

\begin{enumerate}
  \item \textbf{Mô hình hóa.} Tác tử là hai cửa hàng, môi trường là thị trường, trạng thái $(d,c)$, hành động giữ hoặc giảm, thưởng là lợi nhuận. Đây là hệ đa tác tử vì thưởng phụ thuộc hành động đối thủ qua $p_i$ và qua $c_i$.
  \item \textbf{Học và luật.} RL thích nghi được nhưng tốn khám phá và phụ thuộc seed; luật tất định, dễ đọc nhưng cứng. Ở tham chiếu luật đã tối ưu nên RL chỉ có thể hòa.
  \item \textbf{Một bước cập nhật.} Với $Q=\SL{ch8_q3/q_old}$, $r=\SL{ch8_q3/r}$, $\max Q'=\SL{ch8_q3/max_next}$: mục tiêu $\SL{ch8_q3/r}+\SL{ch8_q3/gamma}\times\SL{ch8_q3/max_next}=\SL{ch8_q3/target}$, sai số TD \SL{ch8_q3/td}, giá trị mới \SL{ch8_q3/q_new}, khớp nhật ký chương trình.
  \item \textbf{Khám phá.} Cần ít nhưng không thể bằng 0: $\varepsilon=0$ đạt \SL{ch8_d/khong/eval_mean}, lịch thấp \SL{ch8_d/thap/eval_mean}, lịch cao \SL{ch8_d/cao/eval_mean}.
  \item \textbf{Phần thưởng.} Giảm giá mất 40\% biên lợi nhuận để đổi lấy 15 điểm phần trăm chuyển đổi. Huấn luyện theo doanh thu làm lợi nhuận tụt còn \SL{ch8_g/doanh_thu/eval_mean}: tác tử tối ưu đúng cái được thưởng.
  \item \textbf{Đánh giá.} Phải đóng băng bảng $Q$, tắt khám phá, so cặp trên cùng chuỗi ngẫu nhiên và báo cáo qua nhiều seed.
  \item \textbf{Tốc độ học.} $\alpha$ lớn bám mẫu gần nhất (đổi \SL{ch8_c/0.50/switches} lần), $\alpha$ nhỏ mượt nhưng chậm (\SL{ch8_c/0.01/switches} lần).
  \item \textbf{Phi dừng.} Khi đối thủ cũng học, tỉ lệ giảm giá huấn luyện lên \SL{ch8_h/agg/train_discPct}\% (so với \SL{ch8_e/agg/train_discPct}\%), không seed nào cả hai về \texttt{0000}; môi trường đổi theo thời gian~\cite{ch8:busoniu2008}.
  \item \textbf{Giới hạn.} Bốn trạng thái, không tồn kho, và quan sát thiếu hành động trước của chính tác tử, nên hậu quả ``giảm hôm nay, đối thủ giảm ngày mai'' bị ẩn. Mục~\ref{sec:ch8-env} đo tác động.
  \item \textbf{Tác tử tốt hơn.} Tác tử dựa trên mô hình ở Mục~\ref{sec:ch8-better} đạt tối ưu ở cả năm kịch bản.
\end{enumerate}

\section{Thay đổi môi trường}\label{sec:ch8-env}

Bảng~\ref{tab:ch8-env} và Hình~\ref{fig:ch8-delta} đổi mức giảm $\delta$, tính bền của nhu cầu và cách thưởng. Ở $\delta=5\%$ RL vượt luật, khoảng tin cậy $[\SL{ch8_f/005/ci_lo};\ \SL{ch8_f/005/ci_hi}]$. Ở
$\delta=10\%$ RL lại thua rõ $[\SL{ch8_f/010/ci_lo};\ \SL{ch8_f/010/ci_hi}]$: nó giảm giá khi nhu cầu thấp vì lợi trước mắt mà không thấy phản ứng của đối thủ ở bước sau. ``RL hơn luật'' vì vậy không phải mệnh đề chung.

\bang{ch8_env}{Các biến thể môi trường (10 seed; đánh giá luôn theo lợi nhuận).}{ch8-env}
\hinhpdf{ch8_delta}{1}{Theo mức giảm $\delta$: thưởng của RL, của luật và mức tối ưu; tỉ lệ giảm giá của RL khi đánh giá.}{ch8-delta}

\section{Tác tử cải tiến}\label{sec:ch8-better}

Hai đề xuất nhắm vào lỗ hổng quan sát: Q-learning mở rộng thêm $a_i(t-1)$ vào trạng thái; tác tử dựa trên mô hình ước lượng phản ứng đối thủ, thưởng và nhu cầu từ kinh nghiệm rồi lập kế hoạch đúng 20 bước
(Hình~\ref{fig:ch8-mb}). Tác tử này đạt tối ưu ở \SL{ch8_q10/giam_5/model/n_opt}/10 seed với $\delta=5\%$, trong khi Q-learning gốc chỉ đạt \SL{ch8_q10/giam_5/q/n_opt}/10 (Bảng~\ref{tab:ch8-q10}, Hình~\ref{fig:ch8-regret}). Đổi lại, nó phải biết
cấu trúc trạng thái chung và giả định đối thủ dừng.

\hinhtikz{ch8_mb}{Tác tử dựa trên mô hình: thu kinh nghiệm, cập nhật ba mô hình đếm tần suất, quy hoạch động ngược, chọn hành động.}{ch8-mb}

\bang{ch8_q10}{Ba tác tử trong năm kịch bản.}{ch8-q10}
\hinhpdf{ch8_regret}{1}{Regret mỗi bước của ba tác tử (trung bình mười seed; thấp hơn là tốt hơn).}{ch8-regret}

\FloatBarrier
\begin{luuy}{Phạm vi kết luận}
Các kết quả chỉ có giá trị trong mô hình mô phỏng này, không phải kết luận kinh tế.
\end{luuy}
